% Options for packages loaded elsewhere
% Options for packages loaded elsewhere
\PassOptionsToPackage{unicode}{hyperref}
\PassOptionsToPackage{hyphens}{url}
\PassOptionsToPackage{dvipsnames,svgnames,x11names}{xcolor}
%
\documentclass[
  letterpaper,
  DIV=11,
  numbers=noendperiod]{scrreprt}
\usepackage{xcolor}
\usepackage{amsmath,amssymb}
\setcounter{secnumdepth}{-\maxdimen} % remove section numbering
\usepackage{iftex}
\ifPDFTeX
  \usepackage[T1]{fontenc}
  \usepackage[utf8]{inputenc}
  \usepackage{textcomp} % provide euro and other symbols
\else % if luatex or xetex
  \usepackage{unicode-math} % this also loads fontspec
  \defaultfontfeatures{Scale=MatchLowercase}
  \defaultfontfeatures[\rmfamily]{Ligatures=TeX,Scale=1}
\fi
\usepackage{lmodern}
\ifPDFTeX\else
  % xetex/luatex font selection
\fi
% Use upquote if available, for straight quotes in verbatim environments
\IfFileExists{upquote.sty}{\usepackage{upquote}}{}
\IfFileExists{microtype.sty}{% use microtype if available
  \usepackage[]{microtype}
  \UseMicrotypeSet[protrusion]{basicmath} % disable protrusion for tt fonts
}{}
\makeatletter
\@ifundefined{KOMAClassName}{% if non-KOMA class
  \IfFileExists{parskip.sty}{%
    \usepackage{parskip}
  }{% else
    \setlength{\parindent}{0pt}
    \setlength{\parskip}{6pt plus 2pt minus 1pt}}
}{% if KOMA class
  \KOMAoptions{parskip=half}}
\makeatother
% Make \paragraph and \subparagraph free-standing
\makeatletter
\ifx\paragraph\undefined\else
  \let\oldparagraph\paragraph
  \renewcommand{\paragraph}{
    \@ifstar
      \xxxParagraphStar
      \xxxParagraphNoStar
  }
  \newcommand{\xxxParagraphStar}[1]{\oldparagraph*{#1}\mbox{}}
  \newcommand{\xxxParagraphNoStar}[1]{\oldparagraph{#1}\mbox{}}
\fi
\ifx\subparagraph\undefined\else
  \let\oldsubparagraph\subparagraph
  \renewcommand{\subparagraph}{
    \@ifstar
      \xxxSubParagraphStar
      \xxxSubParagraphNoStar
  }
  \newcommand{\xxxSubParagraphStar}[1]{\oldsubparagraph*{#1}\mbox{}}
  \newcommand{\xxxSubParagraphNoStar}[1]{\oldsubparagraph{#1}\mbox{}}
\fi
\makeatother


\usepackage{longtable,booktabs,array}
\usepackage{calc} % for calculating minipage widths
% Correct order of tables after \paragraph or \subparagraph
\usepackage{etoolbox}
\makeatletter
\patchcmd\longtable{\par}{\if@noskipsec\mbox{}\fi\par}{}{}
\makeatother
% Allow footnotes in longtable head/foot
\IfFileExists{footnotehyper.sty}{\usepackage{footnotehyper}}{\usepackage{footnote}}
\makesavenoteenv{longtable}
\usepackage{graphicx}
\makeatletter
\newsavebox\pandoc@box
\newcommand*\pandocbounded[1]{% scales image to fit in text height/width
  \sbox\pandoc@box{#1}%
  \Gscale@div\@tempa{\textheight}{\dimexpr\ht\pandoc@box+\dp\pandoc@box\relax}%
  \Gscale@div\@tempb{\linewidth}{\wd\pandoc@box}%
  \ifdim\@tempb\p@<\@tempa\p@\let\@tempa\@tempb\fi% select the smaller of both
  \ifdim\@tempa\p@<\p@\scalebox{\@tempa}{\usebox\pandoc@box}%
  \else\usebox{\pandoc@box}%
  \fi%
}
% Set default figure placement to htbp
\def\fps@figure{htbp}
\makeatother

\ifLuaTeX
  \usepackage{luacolor}
  \usepackage[soul]{lua-ul}
\else
  \usepackage{soul}
\fi




\setlength{\emergencystretch}{3em} % prevent overfull lines

\providecommand{\tightlist}{%
  \setlength{\itemsep}{0pt}\setlength{\parskip}{0pt}}



 


\KOMAoption{captions}{tableheading}
\makeatletter
\@ifpackageloaded{bookmark}{}{\usepackage{bookmark}}
\makeatother
\makeatletter
\@ifpackageloaded{caption}{}{\usepackage{caption}}
\AtBeginDocument{%
\ifdefined\contentsname
  \renewcommand*\contentsname{Table of contents}
\else
  \newcommand\contentsname{Table of contents}
\fi
\ifdefined\listfigurename
  \renewcommand*\listfigurename{List of Figures}
\else
  \newcommand\listfigurename{List of Figures}
\fi
\ifdefined\listtablename
  \renewcommand*\listtablename{List of Tables}
\else
  \newcommand\listtablename{List of Tables}
\fi
\ifdefined\figurename
  \renewcommand*\figurename{Figure}
\else
  \newcommand\figurename{Figure}
\fi
\ifdefined\tablename
  \renewcommand*\tablename{Table}
\else
  \newcommand\tablename{Table}
\fi
}
\@ifpackageloaded{float}{}{\usepackage{float}}
\floatstyle{ruled}
\@ifundefined{c@chapter}{\newfloat{codelisting}{h}{lop}}{\newfloat{codelisting}{h}{lop}[chapter]}
\floatname{codelisting}{Listing}
\newcommand*\listoflistings{\listof{codelisting}{List of Listings}}
\makeatother
\makeatletter
\makeatother
\makeatletter
\@ifpackageloaded{caption}{}{\usepackage{caption}}
\@ifpackageloaded{subcaption}{}{\usepackage{subcaption}}
\makeatother
\usepackage{bookmark}
\IfFileExists{xurl.sty}{\usepackage{xurl}}{} % add URL line breaks if available
\urlstyle{same}
\hypersetup{
  pdftitle={HOLPA Key Performance and Agroecological Indicator Protocols},
  pdfauthor={Sarah K. Jones,a,* Andrea Cecilia Sánchez,a,b,c, Chris Dickensd, Matthias S. Gecke, Chaturangi Wickramaratned, Veronique Alaryf,g, Peter Boloh, Dennis Junior Chorumai, Somphasith Douangsavanhj, Modou Gueye Fallk, Gatien Falconnierl,m,n,o, Shweta Guptap,q, Chris Kettler, Smitha Krishnans, Sylvia Nyawirah, Guillermo Orjuela-Ramirezt, Boko Michel Orounladjiu, Piedad Parejav, Telma Sibandal, Christine Lamannae,1},
  colorlinks=true,
  linkcolor={blue},
  filecolor={Maroon},
  citecolor={Blue},
  urlcolor={Blue},
  pdfcreator={LaTeX via pandoc}}


\title{HOLPA Key Performance and Agroecological Indicator Protocols}
\author{Sarah K. Jones,a,* Andrea Cecilia Sánchez,a,b,c, Chris Dickensd,
Matthias S. Gecke, Chaturangi Wickramaratned, Veronique Alaryf,g, Peter
Boloh, Dennis Junior Chorumai, Somphasith Douangsavanhj, Modou Gueye
Fallk, Gatien Falconnierl,m,n,o, Shweta Guptap,q, Chris Kettler, Smitha
Krishnans, Sylvia Nyawirah, Guillermo Orjuela-Ramirezt, Boko Michel
Orounladjiu, Piedad Parejav, Telma Sibandal, Christine Lamannae,1}
\date{2025-12-05}
\begin{document}
\maketitle

\renewcommand*\contentsname{Table of contents}
{
\hypersetup{linkcolor=}
\setcounter{tocdepth}{2}
\tableofcontents
}

\bookmarksetup{startatroot}

\chapter*{Welcome}\label{welcome}
\addcontentsline{toc}{chapter}{Welcome}

\markboth{Welcome}{Welcome}

This is the homepage for the HOLPA Key Performance and Agroecological
Indicator Protocols.

\bookmarksetup{startatroot}

\chapter{Introduction}\label{introduction}

This chapter introduces the HOLPA protocols.

\part{1. Environmental KPIs}

\chapter{1.1 Crop, livestock and fish
diversity}\label{crop-livestock-and-fish-diversity}

\begin{longtable}[]{@{}
  >{\raggedright\arraybackslash}p{(\linewidth - 2\tabcolsep) * \real{0.2222}}
  >{\raggedright\arraybackslash}p{(\linewidth - 2\tabcolsep) * \real{0.7639}}@{}}
\caption{Table 7: Method for calculating key performance
indicators}\tabularnewline
\toprule\noalign{}
\begin{minipage}[b]{\linewidth}\raggedright
\textbf{Indicator}
\end{minipage} & \begin{minipage}[b]{\linewidth}\raggedright
\textbf{Crop, livestock and fish diversity}
\end{minipage} \\
\midrule\noalign{}
\endfirsthead
\toprule\noalign{}
\begin{minipage}[b]{\linewidth}\raggedright
\textbf{Indicator}
\end{minipage} & \begin{minipage}[b]{\linewidth}\raggedright
\textbf{Crop, livestock and fish diversity}
\end{minipage} \\
\midrule\noalign{}
\endhead
\bottomrule\noalign{}
\endlastfoot
Theme & Agrobiodiversity \\
Purpose of indicator & Relevant for the post-2020 Global Biodiversity
Framework Target 10 on `Proportion of agricultural area under productive
and sustainable agriculture', where the Agrobiodiversity Index is a
complementary indicator (CBD/COP/15/L.26). The Agrobiodiversity Index
includes various measures of agrobiodiversity, including crop and
livestock diversity. \\
Data source & Household survey \\
Measurement & Diversity of crop varieties (or species) and fish and
livestock breeds (or species) on the farm, measured using data on
richness and area coverage. \\
Measurement units & Count with associated production area recorded in
hectares, disaggregated by crop and animal species and ideally also
variety/breed \\
Guidance on measurement & \begin{minipage}[t]{\linewidth}\raggedright
This indicator can be calculated by collecting data on the:

\begin{itemize}
\tightlist
\item
  Count of crop varieties and species, and the relative area occupied.
\item
  Count of livestock breeds and species.
\item
  Count of fish breeds and species.
\end{itemize}

These data can be collected as part of a household survey. To verify the
data, and if time allows, the researcher and farmer could conduct a
transect of the farm and discuss and record where different crops are
planted.
\end{minipage} \\
Guidance on data entry and reporting & Be sure to use the same area
units for all species/varieties/breeds, so that diversity indices can be
calculated. Crop and livestock varieties/breeds and species names should
be stored in the local language and also using their scientific and
English names to enable comparisons across sites. A local agronomist may
need to help identify the scientific and English names. \\
Indicator i nterpretation and threshold setting & \\
Limitations & \\
References & FAO (2019), Jones et al.~(2021) \textsuperscript{1,2} \\
\end{longtable}

\chapter{1.2 Insect, mammal and tree
diversity}\label{insect-mammal-and-tree-diversity}

\begin{longtable}[]{@{}
  >{\raggedright\arraybackslash}p{(\linewidth - 2\tabcolsep) * \real{0.2222}}
  >{\raggedright\arraybackslash}p{(\linewidth - 2\tabcolsep) * \real{0.7639}}@{}}
\caption{Table 8: Method for converting key performance indicators from
their original formats to a standardized 0 to 100 scale}\tabularnewline
\toprule\noalign{}
\begin{minipage}[b]{\linewidth}\raggedright
\textbf{Indicator}
\end{minipage} & \begin{minipage}[b]{\linewidth}\raggedright
\textbf{Insect, mammal and tree diversity}
\end{minipage} \\
\midrule\noalign{}
\endfirsthead
\toprule\noalign{}
\begin{minipage}[b]{\linewidth}\raggedright
\textbf{Indicator}
\end{minipage} & \begin{minipage}[b]{\linewidth}\raggedright
\textbf{Insect, mammal and tree diversity}
\end{minipage} \\
\midrule\noalign{}
\endhead
\bottomrule\noalign{}
\endlastfoot
Theme & Biodiversity \\
Purpose of indicator & Relevant for global climate mitigation targets
(the Paris Agreement) and the post-2020 Global Biodiversity Framework
Target 10 on `Proportion of agricultural area under productive and
sustainable agriculture'. Biodiversity indicators selected to be easy to
use and scalable, in line with Essential Biodiversity Variables 2020
initiative recommendations. \\
Data source & Household survey, field survey \\
Measurement & Diversity of insect, mammal and tree species per hectare
of farmland, measured using data on richness, abundance, and density,
and standardised to a scale of 0 to 100 to enable data aggregation. \\
Measurement units & Count with associated production area recorded in
hectares, disaggregated by insect/mammal/tree species, tree maturity,
and local/exotic origin. \\
Guidance on measurement & \begin{minipage}[t]{\linewidth}\raggedright
This indicator can be calculated by collecting data on the:

\begin{itemize}
\tightlist
\item
  Richness and abundance of mammals, insects and trees
\end{itemize}

These data can be collected as part of household and field surveys.

To enable more indepth tree data analysis, e.g.~to support efforts to
monitor native species, ancient specimens, and carbon sequestration
on-farm, additional data can be collected on:

\begin{itemize}
\item
  Proportion of native tree species
\item
  Proportion of mature trees (\textgreater50cm diameter at breast
  height, or height \textgreater10m)
\end{itemize}
\end{minipage} \\
Guidance on data entry and reporting & Species names (if available)
should be stored in the local language and also using their scientific
and English names to enable comparisons across sites. \\
Indicator i nterpretation and threshold setting & \\
Limitations & \\
References & Landmann et al.~(2023)\textsuperscript{3}, Navarro et
al.~(2017)\textsuperscript{4}, GEO-BON\textsuperscript{5} \\
\end{longtable}

\chapter{1.3 Landscape complexity}\label{landscape-complexity}

\begin{longtable}[]{@{}
  >{\raggedright\arraybackslash}p{(\linewidth - 2\tabcolsep) * \real{0.2222}}
  >{\raggedright\arraybackslash}p{(\linewidth - 2\tabcolsep) * \real{0.7639}}@{}}
\toprule\noalign{}
\begin{minipage}[b]{\linewidth}\raggedright
\textbf{Indicator}
\end{minipage} & \begin{minipage}[b]{\linewidth}\raggedright
\textbf{Landscape complexity}
\end{minipage} \\
\midrule\noalign{}
\endhead
\bottomrule\noalign{}
\endlastfoot
Theme & Biodiversity \\
Purpose of indicator & Relevant for the post-2020 Global Biodiversity
Framework Goal B on conserving services provided by ecosystems and
Target 10 on `Proportion of agricultural area under productive and
sustainable agriculture' (CBD/COP/15/L.26). Aligns with SDG indicators
2.4.1 and 15.3.1 on proportion of agricultural area under productive and
sustainable agriculture. Maintaining at least 20\% natural and
semi-natural habitat in agricultural landscapes is important for
safeguarding food security, biodiversity, and ecosystem services
(Garibaldi et al.~2020). \\
Data source & Satellite data, or direct measurement by walking around
natural or semi-natural features on-farm with a GPS. In both cases, the
field or farm boundary needs delineating. \\
Measurement & Percentage of farmland area that is covered with natural
or semi-natural vegetation. \\
Measurement units & Percentage cover, disaggregated by type of
feature \\
Guidance on measurement & \begin{minipage}[t]{\linewidth}\raggedright
We discuss three options for measurement, i) farmer estimates validated
through field surveys (default option in HOLPA), ii) using satellite
imagery, iii) direct measurement.

\textbf{\emph{Farmer estimates}}

A simplistic and rapid way to assess this indicator is to ask
respondents to estimate the proportion of their farmland covered by
natural or semi-natural vegetation (e.g.~bushland, fallow land,
hedgerows, natural grassland, ponds, lakes, forest patches). This can be
validated by assessing the proportion of natural vegetation contained
within three 10x10m plots sampled on the farm, taking a GPS point to
allow for satellite-based analysis of the proportion of natural
vegetation at larger spatial scales.

\textbf{\emph{Satellite imagery}}

Natural or semi-natural features on farmland (see Definitions) are often
visible on high resolution (\textless50m and ideally \textless10m)
satellite imagery, such as Google Earth imagery. In a GIS, import the
boundary of the field or farm of interest (see Definitions). If no
boundary is available, generate a 1km diameter circular plot around each
fieldwork plot, using a GPS point collected at the centre of the plot on
the ground. In Google Earth Engine, digitise all visible natural or
semi-natural features within the area of interest. Calculate the area of
land with natural or semi-natural features to obtain an estimate of the
share of land covered with natural or semi-natural features. Where
possible this should be verified on the ground.

Examples:

\begin{figure}[H]

{\centering \includegraphics[width=4.38542in,height=3.28125in]{kpis/environmental/media/image1.png}

}

\caption{A computer screen shot of a field Description automatically
gene rated}

\end{figure}%

\begin{figure}[H]

{\centering \includegraphics[width=4.42708in,height=3.28125in]{kpis/environmental/media/image2.png}

}

\caption{A screenshot of a computer Description automatically gene
rated}

\end{figure}%

\emph{Known problems:}

\begin{itemize}
\item
  Distinguishing between tree canopies and shade can be difficult.
\item
  Determining what is farmland and what is natural land can be
  difficult.
\item
  Distinguishing between grass and shrubs is not always easy.
\item
  The condition of the vegetation is an unknown but important variable.
\end{itemize}

\textbf{\emph{Direct measurement}}

Natural or semi-natural features coverage can be determined by walking
around the field or farm of interest including all natural or
semi-natural features, recording the boundaries with a handheld GPS.
These areas can be imported into a GIS and used to obtain an estimate of
the percentage of farmland covered with natural or semi-natural
features.

\textbf{Definitions:}

\textbf{\emph{Farmland}} includes areas delineated with generally linear
boundaries and inside which there are one of the following:

\begin{enumerate}
\def\labelenumi{\arabic{enumi}.}
\item
  Regular rows of woody or herbaceous crops
\item
  Visible tractor tracks
\item
  Visible rows of tilled soil
\item
  Evidence of irrigation, e.g.~area is much greener than surrounding
  land, or visible sprinkler system
\item
  Hay bales in field
\end{enumerate}

\textbf{\emph{Natural or semi-natural features}} on farmland includes
linear or areal features on farmland comprising:

\begin{itemize}
\item
  Single or small groups of trees or shrubs not regularly arranged
  within the field. These include woodlots, forest remnants.
\item
  Narrow strips of trees, shrubs or other vegetation along field edge
  that have no visible management, e.g., are not arranged in rows, no
  signs of tractor or tillage activity, irregular shape. These include
  hedgerows, flower strips, grass borders.
\item
  Other natural or semi-natural features, including long-term
  (\textgreater6 month) fallow land, rivers, ponds.
\end{itemize}

This list of natural or semi-natural features should be discussed and
adjusted with local communities and/or based on expert knowledge of the
context
\end{minipage} \\
Guidance on data entry and reporting & Record which natural or
semi-natural features were detected, e.g.~hedgerows, flower strip. For
fallow land, record how long the land has been in fallow. \\
Indicator i nterpretation and threshold setting & \\
Limitations & \\
References & Garibaldi et al.~(2020), Mohamed et
al.~(2024)\textsuperscript{5,6} \\
\end{longtable}

\chapter{1.4 Water stress}\label{water-stress}

\begin{longtable}[]{@{}
  >{\raggedright\arraybackslash}p{(\linewidth - 2\tabcolsep) * \real{0.2222}}
  >{\raggedright\arraybackslash}p{(\linewidth - 2\tabcolsep) * \real{0.7639}}@{}}
\toprule\noalign{}
\begin{minipage}[b]{\linewidth}\raggedright
\textbf{Indicator}
\end{minipage} & \begin{minipage}[b]{\linewidth}\raggedright
\textbf{Water stress}
\end{minipage} \\
\midrule\noalign{}
\endhead
\bottomrule\noalign{}
\endlastfoot
Theme & Water \\
Purpose of indicator & Aligns to SDG 6.4.2 on level of water stress, and
SDG 6.6.1 on extent of water-related ecosystems. \\
Data source & Household survey \\
Measurement & Proportion of months with agricultural water shortages per
year, for all agriculturally active months \\
Measurement units & \% \\
Guidance on measurement & \begin{minipage}[t]{\linewidth}\raggedright
This indicator can be calculated by collecting data on the:

\begin{itemize}
\tightlist
\item
  Number of months for which a household experienced water shortages in
  the last 12 months, in relation to crop, livestock, fish or other
  agricultural activities (ideally disaggregated).
\end{itemize}

The indicator is calculated by summing the number of months of water
stress and dividing this by 12, and then taking the inverse (1/x) so
that higher values indicate less water stress.

Additional data can be collected on water sources (disaggregated by
rainwater, water storage on-farm, water storage off-farm, flood water,
piped water, other) to distinguish rainfed water stress versus irrigated
water stress, and whether water harvesting measures are used.

These data can be collected as part of a household survey.
\end{minipage} \\
Guidance on data entry and reporting & Collect data by asking which
months (Jan, Feb, Mar, etc) water is needed per activity, and then which
months there are water shortages per activity, to ensure alignment
between the two responses. \\
Indicator i nterpretation and threshold setting & This indicator
presents the percentage of months in a year with water stress. Lower
values indicate less stress, while high values indicate more stress.
Therefore the indicator values should be inversed \\
Limitations & \begin{minipage}[t]{\linewidth}\raggedright
Key limitations of the measurement approach for this indicator include:

\begin{itemize}
\item
  Water quality is not accounted for.
\item
  Water storage and delivery infrastructure are not considered.
\end{itemize}

Alternatives to consider for future surveys include the
\textquotesingle\textquotesingle Agricultural Water Scarcity
Index\textquotesingle\textquotesingle{} which is a dimensionless index;
or \textquotesingle\textquotesingle Water
Productivity\textquotesingle\textquotesingle{} (kg produced/m3 water
consumed); or Quantify freshwater (rainfall, groundwater and surface
water) supplied to different agricultural enterprises on the farm during
the month or year. For the latter, users should:

\begin{itemize}
\item
  Estimate amount of water (m3/year) required by crop, livestock, fish
  or other agricultural enterprises (disaggregated) on the farm.
\item
  Estimate the amount of water (m3/year) supplied to each agricultural
  enterprise on the farm.
\item
  Calculate: Water Stress (\%) = (Total freshwater supplied \textbf{/}
  Available freshwater resources) ×

  \begin{enumerate}
  \def\labelenumi{\arabic{enumi}.}
  \setcounter{enumi}{99}
  \tightlist
  \item
  \end{enumerate}
\item
  Classify water stress (\%) into the following categories:

  \begin{itemize}
  \item
    No stress = \textless10. Low water demand relative to water supply.
  \item
    Low stress = 10-20. Moderate demand -- sustainable status.
  \item
    Medium stress = 20-40. Significant pressure on water resources.
  \item
    High stress = 40-80. Strong competition for water -- management is
    critical.
  \item
    Critical stress = \textgreater80. Unsustainable water extraction for
    agriculture.
  \end{itemize}
\end{itemize}
\end{minipage} \\
{[}Note to me: add these to zotero{]}\{. comment-start id=``17''
author=``J ones, Sarah ( Alliance Biov ersity-CIAT)'' d ate=``2025-12-
05T11:56:00Z'' \}References{[}{]} \{.comment-end id=``17''\} &
\begin{minipage}[t]{\linewidth}\raggedright
\begin{enumerate}
\def\labelenumi{\arabic{enumi}.}
\item
  Biancalani, R., \& Marinelli, M. (2021). Assessing SDG indicator 6.4.2
  "level of water stress" at major basins level. \emph{UCL Open
  Environment}.
\item
  Food and Agriculture Organization (FAO). (2021). \emph{SDG Indicator
  6.4.2 -- Level of water stress: Methodology}. United Nations FAO.
  (\url{http://www.fao.org/aquastat}).
\item
  van Vliet, M. T. H., Jones, E. R., et al. (2021). Global water
  scarcity including surface water quality and expansions of clean water
  technologies. \emph{Environmental Research Letters}, 16(2), 024020.
  \url{https://doi.org/10.1088/1748-9326/ABBFC3}
\end{enumerate}
\end{minipage} \\
\end{longtable}

\chapter{\texorpdfstring{1.5 \hl{Energy
use}}{1.5 Energy use}}\label{energy-use}

\begin{longtable}[]{@{}
  >{\raggedright\arraybackslash}p{(\linewidth - 2\tabcolsep) * \real{0.2222}}
  >{\raggedright\arraybackslash}p{(\linewidth - 2\tabcolsep) * \real{0.7639}}@{}}
\toprule\noalign{}
\begin{minipage}[b]{\linewidth}\raggedright
\textbf{Indicator}
\end{minipage} & \begin{minipage}[b]{\linewidth}\raggedright
{[}Sustainable energy security aims to meet present energy needs without
compromising the ability of future generations to meet their own needs.
This requires a transition to cleaner energy sources, improved energy
efficiency, and the development of more sustainable patterns of
consumption. (https://e
nergy.sustainability-directory.com/question/what-are
-key-indicators-of-energy-security/){]}\{.comment-start id=``19'' author
=``CHIMONYO, Vimbayi Grace Petrova (CIMMYT-Zimbabwe)''
date=``2025-11-26T10:21:00Z''\}{[}Renewable Energy Share -Percentage of
energy supply from renewable sources{]}\{.comment-start id=``20'' author
=``CHIMONYO, Vimbayi Grace Petrova (CIMMYT-Zimbabwe)''
date=``2025-11-26T10:32:00Z''\}\phantomsection\label{21}{This will be
hard to measure-\/-how do we define energy share without measuring hours
of use, energy efficiency level etc, and some estimate of KJ used per
source? For example: one hour of LPG use is very different from 1 hour
of bio-mass burning stoves.}\textbf{Sustainable energy
use}\phantomsection\label{19}{\phantomsection\label{20}{\phantomsection\label{21}{}}}
\end{minipage} \\
\midrule\noalign{}
\endhead
\bottomrule\noalign{}
\endlastfoot
Theme & Energy use \\
Purpose of indicator & {[}we need to consider that most of the global
south is affected by energy poverty and are not the main consumers of
fossil fuels for energy. another point to consider some energy sources
which can be considered green energy - e.g.~biogas digestors contribute
to GHG emissions when they leak.{]}\{.comment-start id=``22'' author
=``CHIMONYO, Vimbayi Grace Petrova (CIMMYT-Zimbabwe)''
date=``2025-11-26T09:24:00Z''\}Fossil fuels and other non-renewable
energy sources\phantomsection\label{22}{} contribute to global
greenhouse gas emissions, exacerbating global warming. Increasing the
share of renewable energy sources while reducing overall energy
consumption contributes to net zero emissions targets (Paris Agreement,
SDG 13), while self- and locally produced energy helps ensure an
affordable and reliable energy supply (SDG 7). \\
Data source & Household survey \\
\begin{minipage}[t]{\linewidth}\raggedright
{[}Can we define this as a binary variable\\
1- use only clean energy sources for domestic use (LPG/Electric for
cooking, e lectric/solar for heating and lighting)\\
0-use unclean or mix of clean and unclean sources\\
Likewise for on-farm use-\/-1 if using electric, solar or manual tools
/equipment{]}\{. comment-start id=``23'' author='' Alvi, Muzna (
IFPRI-India)'' da te=``2025-12-0 2T14:56:00Z''\} Measurement{[}{]}
\{.comment-end id=``23''\}\strut
\end{minipage} & \begin{minipage}[t]{\linewidth}\raggedright
\begin{itemize}
\item
  Proportion of energy sources that are from renewable sources (energy
  produced from solar panels, wind turbines, hydropower or other
  renewable sources).
\item
  Proportion of consumed energy that is produced on-farm or by local
  suppliers.
\end{itemize}
\end{minipage} \\
Measurement units & \% \\
Guidance on measurement & \begin{minipage}[t]{\linewidth}\raggedright
This indicator can be calculated by collecting data on:

\begin{itemize}
\tightlist
\item
  Sources of energy , disaggregated by renewable and non-renewable, per
  activity (cooking, irrigation, tillage, and other
  food-related/agricultural activities){[}. Used to provide a Likert
  scale measure of clean energy:{]}\{.insertion author=``Jones, Sarah
  (Alliance Bioversity-CIAT)''
\end{itemize}

date=``2025-12-05T10:14:00Z''\}{}

\begin{verbatim}
-   [Totall]{.insertion

author="Jones, Sarah (Alliance Bioversity-CIAT)"
    date="2025-12-05T10:14:00Z"}[y]{.insertion

author="Jones, Sarah (Alliance Bioversity-CIAT)"
    date="2025-12-05T10:15:00Z"} [clean: 100%
    renewable energy sources]{.insertion

author="Jones, Sarah (Alliance Bioversity-CIAT)"
\end{verbatim}

date=``2025-12-05T10:14:00Z''\}{[}{]}\{.paragraph-insertion

\begin{verbatim}
author="Jones, Sarah (Alliance Bioversity-CIAT)"
    date="2025-12-05T10:14:00Z"}

-   [Partially clean:]{.insertion

author="Jones, Sarah (Alliance Bioversity-CIAT)"
    date="2025-12-05T10:14:00Z"}
    [1-99]{.insertion

author="Jones, Sarah (Alliance Bioversity-CIAT)"
    date="2025-12-05T10:15:00Z"}[%]{.insertion

author="Jones, Sarah (Alliance Bioversity-CIAT)"
\end{verbatim}

date=``2025-12-05T10:14:00Z''\}{[}{]}\{.paragraph-insertion

\begin{verbatim}
author="Jones, Sarah (Alliance Bioversity-CIAT)"
    date="2025-12-05T10:14:00Z"}

-   [Not clean: 0%]{.insertion

author="Jones, Sarah (Alliance Bioversity-CIAT)"
    date="2025-12-05T10:15:00Z"}
\end{verbatim}

\begin{itemize}
\tightlist
\item
  Proportion of energy used that is produced on-farm, by the community,
  or purchased from other sources
\end{itemize}

Data on the proportion of renewable and locally produced energy sources
can be combined to provide a Likert-scale measure of the sustainability
of energy use.
\end{minipage} \\
Guidance on data entry and reporting & {[}\hl{Record the energy sources
in the system assessed.}{]}\{.insertion author =``CHIMONYO, Vimbayi
Grace Petrova (CIMMYT-Zimbabwe)'' date=``2025-11-26T07:33:00Z''\} \\
Indicator i nterpretation and threshold setting & \\
Limitations & \\
References & {[}\hl{https://ener
gy.sustainability-directory.com/question/what-are-ke
y-indicators-of-energy-security/}{]}\{.insertion author =``CHIMONYO,
Vimbayi Grace Petrova (CIMMYT-Zimbabwe)''
date=``2025-11-26T08:22:00Z''\} \\
\end{longtable}

\chapter{1.6 Climate mitigation}\label{climate-mitigation}

\begin{longtable}[]{@{}
  >{\raggedright\arraybackslash}p{(\linewidth - 2\tabcolsep) * \real{0.1528}}
  >{\raggedright\arraybackslash}p{(\linewidth - 2\tabcolsep) * \real{0.8333}}@{}}
\toprule\noalign{}
\begin{minipage}[b]{\linewidth}\raggedright
\textbf{Ind icator}
\end{minipage} & \begin{minipage}[b]{\linewidth}\raggedright
\textbf{Net greenhouse gas emissions}
\end{minipage} \\
\midrule\noalign{}
\endhead
\bottomrule\noalign{}
\endlastfoot
Theme & Climate mitigation \\
Purpose of i ndicator & Globally, the agriculture and land use sector
contributes approximately 18\% of total greenhouse gas emissions, rising
to nearly 60\% of GHG emissions in parts of the developing world where
agriculture dominates the economy. With the world set to exceed the
emissions threshold for keeping warming below 1.5\textsuperscript{o}C in
less than 10 years\textsuperscript{7}, mitigation from all sectors is a
top global priority.

This indicator is relevant for global climate mitigation goals (Paris
Agreement) and nationally determined contributions (NDCs). It is also
relevant for the post-2020 Global Biodiversity Framework Target 10 on
`Proportion of agricultural area under productive and sustainable
agriculture', where changes in soil organic carbon stocks is a
complementary indicator (CBD/COP/15/L.26). \\
Data source & Farmer survey, literature review, and/or field
measurements \\
Mea surement & Qualitative score of net greenhouse gas emissions

(alternative: carbon storage in soil and above-ground biomass) \\
Mea surement units & 5-point Likert scale (alternative: t C
ha\textsuperscript{−1} yr\textsuperscript{−1}) \\
Guidance on mea surement & \begin{minipage}[t]{\linewidth}\raggedright
We discuss two options for measurement of net greenhouse gas emissions:
i) qualitative assessment of climate mitigation potential of on-farm
practices, ii) quantitative assessment of carbon storage and
sequestration only, based on soil organic carbon (SOC) and estimates of
carbon stored in live woody biomass.

\textbf{Qualitative assessment}

Given the importance of climate change mitigation and the relative
difficulty of accurately quantifying net GHG emissions from smallholder
agriculture across diverse country contexts, we propose a qualitative
approach for integration of mitigation into all performance assessments.

Measurement involves:

\begin{itemize}
\item
  Classifying farming practices according to their net emissions:
  1,2,3,4,5, where 1 represents high emitting and low sequestrating
  practices, and 5 represents low emitting and high sequestrating
  practices.
\item
  Collecting household survey data on the share of cropland under each
  practice.
\item
  Calculating a weighted average score for all practices found on the
  farm, using the share of cropland as weights.
\end{itemize}

Note that the final score could instead be weighted by productivity for
an emissions intensity estimate, if practice data are collected for each
main crop/livestock/fish.

Locally relevant farming practices need to be identified and classified
into a net emissions category. The classification can be done through
literature review and local experts.

Some qualitative assessments/reviews of mitigation potential of
agricultural practices already exist. We use one such review by Bell et
al.~(2018) to provide an example scoring of practices by emission
category. The overall score is an average of the component score. In the
overall score, 1 represents high emitting and low sequestrating
practices, and 5 represents low emitting and high sequestrating
practices.

\section{---------------------------------}\label{section}

\textbf{Practices} \textbf{Carbon }Carbon \textbf{Avoided GHG }Overall
\textbf{References} s equestration sequestration emissions**
score\textbf{ i n soils} in biomass\textbf{ ---------------
--------------- - -------------- ------------- -----------
---------------- }*Cropping

systems***

Agroforestry 5 5 4 4.66 Bell et al.

\begin{verbatim}
                                               (2018)
\end{verbatim}

Biochar 3 3 2 2.66 Bell et al.

\begin{verbatim}
                                               (2018)
\end{verbatim}

Crop rotation 3 3 3 3 Bell et al.

\begin{verbatim}
                                               (2018)
\end{verbatim}

Drip irrigation 3 3 2 2.66 Bell et al.

\begin{verbatim}
                                               (2018)
\end{verbatim}

Embedded 4 5 4 4.33 Bell et al. natural (2018) (hedgerows)

Green manure 4 3 2 3 Bell et al.

\begin{verbatim}
                                               (2018)
\end{verbatim}

Intercropping 5 3 5 4.33 Bell et al.

\begin{verbatim}
                                               (2018)
\end{verbatim}

Microdosing 3 3 2 2.66 Bell et al.

\begin{verbatim}
                                               (2018)
\end{verbatim}

Mulching 4 3 3 3.33 Bell et al.

\begin{verbatim}
                                               (2018)
\end{verbatim}

Reduced tillage 4 3 4 3.66 Bell et al.

\begin{verbatim}
                                               (2018)
\end{verbatim}

***Lowland rice

systems***

Biochar 4 4 2 3.33 Bell et al.

\begin{verbatim}
                                               (2018)
\end{verbatim}

Green manure 4 4 4 4 Bell et al.

\begin{verbatim}
                                               (2018)
\end{verbatim}

Microdosing 4 4 2 3.33 Bell et al.

\begin{verbatim}
                                               (2018)
\end{verbatim}

Rice-fish 3 3 4 3.33 Bell et al. integration (2018)

Alternate 3 3 4 3.33 Bell et al. wetting and (2018) drying

***Livestock

systems***

Improved feed 4 4 4 4 Bell et al. quality (2018)

Planting 4 3 3 3.33 Bell et al. N-fixing (2018) legumes

Rotational 4 3 3 3.33 Bell et al. grazing (2018)

***Fish

systems***

Improved feed 3 3 3 3 Bell et al. quality (2018)
---------------------------------
---------------------------------------------------------

\textbf{Quantitative assessment}

Measurement involves collecting data on tree diameter at breast height
(dbh), tree height, number of trees, and collecting data on soil organic
carbon through laboratory analysis of soil samples.

Carbon storage and sequestration in woody biomass:

\begin{itemize}
\item
  Get on-farm tree count for trees with dbh \textgreater{} 50cm or
  \textgreater10m in height
\item
  Convert above and below-ground biomass to t/ha using IPCC conversion
  factors.
\end{itemize}

Soil carbon storage:

\begin{itemize}
\tightlist
\item
  Use the median SOC in the 0-20cm zone from across the three plots per
  farm (see Soil Health indicator).
\end{itemize}

\textbf{Additional GHG measurement options (beyond the scope of the
standard HOLPA tool, see Complementary Indicators)}

The most precise, quantitative measurements of GHG emissions are
obtained through field measurements of GHG fluxes, either through
collection of gas samples and later analysis in a gas chromatography lab
(e.g.~for ruminant methane emission), or through real-time measurement
of gas fluxes using an infrared gas analyzer (e.g.~leaf photosynthesis
and respiration, soil respiration, or landscape eddy covariance). These
methods require costly equipment and/or laboratory analysis, as well as
expert technicians to collect and analyze the samples. Further details
on the procedures are provided in the section on Complementary
Indicators.

In lieu of direct measurement, all other approaches to quantifying GHG
emissions from agriculture rely on collecting activity data and using
emissions factors to convert activities into emissions levels (e.g.~from
IPCC). Activity data is a description of the farm practices and areas
(or number of livestock) over which they are practiced. Activity data
can be collected via household surveys, potentially supplemented with
measurement of farm sizes. Higher levels of detail in the activity data
(amounts of inputs used, crops and crop varieties, soil types, feed
formulations, tree species in agroforestry systems, etc.) generally
allow for more precise estimation of emissions. However, conversion of
activity data into quantitative emissions relies on the use of emissions
factors. Most widely available emissions factors (e.g.~Tier 1 Approach)
are based on data from temperate countries
(\href{https://onlinelibrary.wiley.com/doi/10.1111/gcb.12361}{Ogle et
al. 2013}; Albanito et al.~2017) and may not be appropriate for
estimating GHG emissions from agricultural activities in tropical
countries.

Models of GHG emissions also rely on collection of activity data and
emissions factors, but many can account for local conditions, such as
soil types, livestock breeds, climate, etc to produce more appropriate
emissions estimates for tropical or developing countries not well
represented in the Tier 1 approach. Several models are available for
calculating emissions, including the Ex-Ante Carbon Balance Tool
(\href{https://www.fao.o\%20rg/in-action/epic/ex-act-tool/suite-of-tools/ex-act/en/}{EXACT})
and the \href{https://coolfarmtool.org/}{Cool Farm Tool}, however use of
these models requires a degree of technical expertise.
\end{minipage} \\
Guidance on data entry and r eporting & \\
I ndicator interp retation and t hreshold setting & \\
Lim itations & \\
Re ferences & \begin{minipage}[t]{\linewidth}\raggedright
\begin{itemize}
\item
  Thornton et al.~2018. A qualitative evaluation of CSA options in mixed
  crop-livestock systems in developing countries.\textsuperscript{8}
\item
  Bell et al.~2018. A Practical Guide to Climate Smart Agricultural
  Technologies in Africa. \emph{CCAFS Working Paper
  No.~224}.\textsuperscript{9}
\item
  Anuga et al.~2020. Towards low carbon agriculture: Systematic
  narratives of climate-smart agriculture mitigation potential in
  Africa. \emph{Current Research in Environmental
  Sustainability}.\textsuperscript{10}
\item
  MacLeod et al.~2020. Quantifying greenhouse gas emissions from global
  aquaculture\emph{. Scientific Reports}.\textsuperscript{11}
\end{itemize}
\end{minipage} \\
\end{longtable}

\part{2. Agricultural KPIs}

\chapter{2.1 Crop health}\label{crop-health}

\section{Indicator}\label{indicator}

\textbf{Crop health}

\section{Theme}\label{theme}

Crop health

\section{Purpose of indicator}\label{purpose-of-indicator}

Crop and pasture health can influence agricultural productivity and
resilience of plants to stress, pests and competition. Crop health can
be affected by any biological, chemical, or physical factors that may
affect plant physiology and crop performance, impacting food production
(SDG2) and the spread of food-borne diseases with consequences for human
health (SDG3).

\section{Data source}\label{data-source}

Household survey, field measurements

\section{Measurement}\label{measurement}

\% crop losses, and SOCLA indicators of crop health

\section{Measurement units}\label{measurement-units}

Score from 1 to 10, representing the mean indicator value across 10
indicators

\section{Guidance on measurement}\label{guidance-on-measurement}

We propose two complementary measures of crop health: i) Qualitative
based on household survey responses, ii) Qualitative based on field
observations.

\textbf{Household survey based}

A simple measure of crop health can be provided by collecting responses
on: What percentage of the total crop production was lost or damaged in
the last 12 months?

\textbf{Field based}

More detailed data on crop health can be determined through field
measurements interpreted jointly by farmers and researchers. Field
measurements to collect data on each crop health indicator should be
conducted every 50m along a transect of approximately 150m (this may
vary depending on the farm size), to collect data on all 10 crop health
indicators at three sampling sites per farm. Crop health should be
assessed for only the main crop(s) grown on the farm to keep the task
manageable.

At three 10x10m square plots per farm, observe the crops and/or pasture
present in the plot and jointly with the respondent record the responses
to the indicators with the respondent. Refer to the example images
provided in the field survey (Appendix 7.1.3.3 Crop and pasture health)
to determine the most appropriate responses crop health related
questions.

The proposed indicators were developed by the Latin American Society for
Agroecology (SOCLA) and described in Nicholls et
al.~(2004)\textsuperscript{12}. The SOCLA 10 indicators of crop health
are:

\begin{enumerate}
\def\labelenumi{\arabic{enumi}.}
\tightlist
\item
  Appearance
\item
  Crop growth
\item
  Disease incidence
\item
  Insect pest incidence
\item
  Natural enemy abundance and diversity
\item
  Weed competition and pressure
\item
  Actual or potential yield
\item
  Vegetational diversity
\item
  Natural surrounding vegetation
\item
  Management system
\end{enumerate}

Each indicator is scored from 1 to 10 based on the following guidelines
(from Nicholls et al.~2004, Table 1):

\pandocbounded{\includegraphics[keepaspectratio]{kpis/agricultural/media/image3.png}}

\section{Guidance on data entry and
reporting}\label{guidance-on-data-entry-and-reporting}

Record the crop species and variety in the system assessed. Be sure to
take a GPS point and keep a record of the geographic coordinates where
crop health was assessed to allow follow-up surveys in the same
locations.

\section{Indicator interpretation and threshold
setting}\label{indicator-interpretation-and-threshold-setting}

\section{Limitations}\label{limitations}

\section{References}\label{references}

Nicholls et al.~(2004), FAO (2019) \textsuperscript{13},
\textsuperscript{1}

\chapter{2.2 Animal health}\label{animal-health}

\section{Indicator}\label{indicator-1}

\textbf{Animal health}

\section{Theme}\label{theme-1}

Animal health

\section{Purpose of indicator}\label{purpose-of-indicator-1}

Maintaining healthy and diverse stocks of animals on-farm contributes to
overall agroecosystem functioning for more sustainable farming systems
(SDG2, GBF Target 10). Keeping animals in poor health increases the
spread of zoonoses, food-borne disease, and livestock mortalities, with
associated negative consequences on human health (SDG2 and 3) and
livelihoods (SDG1). Healthier animals have a smaller environmental
footprint (SDG12 and 13).

\section{Data source}\label{data-source-1}

Household survey with farm transect

\section{Measurement}\label{measurement-1}

Qualitative measure of livestock disease and losses. Alternative: Health
and welfare of farm animals, measured using the welfare outcome
assessment tool as applied to each livestock species.

\section{Measurement units}\label{measurement-units-1}

Various (e.g.~number of animals in each welfare category, number of
mortalities per 100 animals)

\section{Guidance on measurement}\label{guidance-on-measurement-1}

We discuss two measures of animal health: i) qualitative welfare
assessment, ii) (optional) full welfare assessment.

\textbf{Qualitative asessment}

Likert score (1 to 4) based on the following questions:

What was the extent of injury, illness, or death of livestock due to
diseases in the last 12 months?

\begin{enumerate}
\def\labelenumi{\arabic{enumi}.}
\tightlist
\item
  None
\item
  Low
\item
  Medium
\item
  High
\end{enumerate}

What was the extent of injury, illness or death of fish due to diseases
in the last 12 months?

\begin{enumerate}
\def\labelenumi{\arabic{enumi}.}
\tightlist
\item
  None
\item
  Low
\item
  Medium
\item
  High
\end{enumerate}

\textbf{Full welfare assessment (optional)}

The welfare outcome assessment tool is a practical and scientifically
informed way of assessing and measuring animal welfare. It aims to
provide an objective, accurate and direct picture of animal welfare.

The assessment measures the impact of farm inputs (e.g.~housing, space,
feed, veterinary care, management) on the health, physical condition and
behaviour of the animals themselves. There are different assessment
approaches for each livestock species: laying hens, dairy cows, pigs,
broiler hens, beef cattle, sheep. The method was developed by Assurewel
and is implemented by the UK Soil Association and RSPCA:
http://www.assurewel.org/aboutassurewel/aboutwelfareoutcomeassessment.html.

The assessment of one livestock species takes around 30 mins per farm.
The researcher will need to complete, with the farmer, the AssureWel
simple scoresheet for each livestock produced on the farm. For example,
for dairy cattle, this includes documenting:

\begin{itemize}
\tightlist
\item
  Mobility (good, impaired, severely impaired)
\item
  Body condition (thin, moderate, fat)
\item
  Cleanliness (clean to very dirty)
\item
  Hair loss, lesions and swellings
\item
  Number of cows with broken tails
\item
  Number of cows needing further care
\item
  Farmer records of mobility, mastitis, calf and heifer survivability,
  cull and casualty cows.
\end{itemize}

\section{Guidance on data entry and
reporting}\label{guidance-on-data-entry-and-reporting-1}

\section{Indicator interpretation and threshold
setting}\label{indicator-interpretation-and-threshold-setting-1}

\section{Limitations}\label{limitations-1}

\section{References}\label{references-1}

\begin{itemize}
\tightlist
\item
  \textsuperscript{14} -- assessment protocol for dairy cattle
\item
  \textsuperscript{15} -- scoresheet for dairy cattle
\item
  \textsuperscript{16} -- explanation of measures for dairy cattle
\item
  See \url{http://www.assurewel.org} for protocols, scoresheets and
  explanations for laying hens, pigs, broiler hens, beef cattle, and
  sheep.
\end{itemize}

\chapter{2.3 Soil health}\label{soil-health}

\section{Purpose of indicator}\label{purpose-of-indicator-2}

Soil organic carbon or organic matter (SOC or SOM) is a key indicator of
soil health, supporting vegetation growth, soil stability, soil
fertility and other vital ecosystem services. SOC/SOM assessment is
relevant for assessing carbon emission and sequestration, and soil
health, under local and international commitments. This indicator is
relevant to global climate mitigation goals (Paris Agreement) and the
post-2020 Global Biodiversity Framework Target 10 on `Proportion of
agricultural area under productive and sustainable agriculture', where
changes in soil organic carbon stocks is a complementary indicator
(CBD/COP/15/L.26). In resource-constrained projects, it can be
challenging to collect soil samples and run laboratory analyses to
calculate SOC. We propose using simple qualitative measures of soil
erosion and fertility o provide an alternative, rapid assessment of soil
health in cases where SOC is unattainable.

\section{Data source}\label{data-source-2}

Soil samples collected on-site (alternative: household survey)

\section{Measurement}\label{measurement-2}

Laboratory measurement of \% soil organic carbon, or \% soil organic
matter (alternative: qualitative assessment of erosion and fertility
through household survey)

\section{Measurement units}\label{measurement-units-2}

\% (alternative: Likert-scale score)

\section{Guidance on measurement}\label{guidance-on-measurement-2}

We discuss two options for the measurement of soil health: i) soil
organic carbon calculated in a laboratory from soil samples, ii)
qualitative measure of soil health.

\textbf{Soil organic carbon}

Soil samples should be taken randomly, in a W-shaped pattern, from 5
points per plot, for three 10x10m plots per farm.

After collecting the soil samples, a standard laboratory procedure
should be used to determine the carbon content of the sample. The
loss-on-ignition method17 can be used to estimate the SOC/SOM using the
following steps.

Air-dry the soil samples after collection, and grind fine enough to
sieve through a mesh \textless2 mm or smaller.

Oven dry each sample at 105°C for 24 hours.

Weigh 5.000g ±0.001g of each oven-dried sample and place each into a
crucible. This will be the pre-ignition weight of the sample.

Heat the sample in a muffle furnace at 375°C for 16 hours or overnight

Weigh the sample after the ignition and this will be the post-ignition
weight

Calculate percentage SOC/SOM using the following equation:

SOC~\left(SOM\right)\left(\%\right)=\frac{pre\ ignition\ weight\ \left(g\right)-post\ ignition\ weight\ \left(g\right)}{pre\ ignition\ weight\ (g)}

Alternatively, a Carbon Nitrogen (CN) Elemental Analyser can be used to
estimate soil organic carbon content. The C-N elemental analyser
provides a quantitative measurement of SOC and N content based on the
amount of C and N produced after combustion of products in a sample.

Alternatively, the Walkey Black method can be used to estimate soil
organic matter content. Walkey Black provides a qualitative measurement
of SOM after oxidising the soil with potassium dichromate and H2SO4 and
measuring the color change; which is proportionate to the organic matter
available in the soil. A conversion factor can be used to convert either
SOC to SOM or vice versa.

\textbf{Qualitative measure of soil health}

To complement SOC data, and to allow for a basic assessment of soil
health in cases where neither soil organic carbon data or SOCLA (see
Complementary indicators) can be collected (e.g.~to due to resource
constraints), HOLPA includes two standard questions on farmer
perceptions of soil health. Farmer perceptions of soil fertility and
erosion are gathered using a simple Likert scale, based on the following
two questions:

How would you describe the fertility of your soil on your farmland?

\begin{verbatim}
- Highly fertile (score 5)
- Moderately fertile (score 3.66)
- Low fertility (score 2.33)
- Infertile (score 1).
\end{verbatim}

How would you describe the level of soil erosion problems on your
farmland?

\begin{verbatim}
- Soil erosion is not a problem on my farm (score 5)
- Soil erosion is a minor problem on my farm (score 3)
- Soil erosion is a major problem on my farm (score 1).
\end{verbatim}

These perception questions can be used to assign a soil health score
from 1 (infertile and major soil erosion) to 3 (highly fertile with no
soil erosion), by taking the median across the fertility and erosion
responses.

\section{Guidance on data entry and
reporting}\label{guidance-on-data-entry-and-reporting-2}

Report the SOC content as a volume-based concentration (i.e.~\%)

\section{Indicator interpretation and threshold
setting}\label{indicator-interpretation-and-threshold-setting-2}

\section{Limitations}\label{limitations-2}

\section{References}\label{references-2}

Bahadori, M., \& Tofighi, H. (2016). A modified Walkley-Black method
based on spectrophotometric procedure. Communications in Soil Science
and Plant Analysis, 47(2), 213-220. 18

Gelman, F., Binstock, R., \& Halicz, L. (2012). Application of the
Walkley--Black titration for the organic carbon quantification in
organic rich sedimentary rocks. Fuel, 96, 608-610.19

Nelson, D.W., and Sommers, L.E. (1996). Total Carbon, Organic Carbon,
and Organic Matter. In Methods of Soil Analysis (John Wiley \& Sons,
Ltd), pp.~961--1010. https://doi.org/10.2136/sssabookser5.3.c34.17

Walkley, A. and I.A. Black. 1934. An examination of the Degtjareff
method for determining soil organic matter and a proposed modification
of the chromic acid titration method. Soil Sci. 37:29-38.20

\chapter{2.4 Nutrient use}\label{nutrient-use}

\section{Indicator}\label{indicator-2}

\textbf{Nutrient use}

\section{Theme}\label{theme-2}

Nutrient management

\section{Purpose of indicator}\label{purpose-of-indicator-3}

Nutrient use per unit area provides information on the under or over-use
of fertilizers, which has implications for crop productivity, economic
viability, and environmental sustainability. By monitoring nutrient
applications, farmers and agronomists can determine if the current
inputs are meeting or exceeding the nutrient needs of the crop and if
adjustments are needed to optimize input use and reduce nutrient losses
to the environment. This indicator is related to a number of Sustainable
Development Goals (SDG) e.g., SDG 2 on zero hunger and SDG 15 life on
land.

\section{Data source}\label{data-source-3}

Field measurements, household survey

\section{Measurement}\label{measurement-3}

Fertilizer applications in the production system

\section{Measurement units}\label{measurement-units-3}

Fertilizer input (kg) per hectare per year, disaggregated by organic and
chemical inputs. Can be presented as a ratio of recommended fertilizer
input intensity in the landscape obtained from extension services and
literature.

\section{Guidance on measurement}\label{guidance-on-measurement-3}

Nutrient use can be determined at the three levels i.e.~plot level (the
level of individual crops), farm level (the level of a whole farm and
its cropping system), and landscape level (a larger area e.g.~region or
watershed).

At all scales, it is important to consider all possible sources of
inputs into the system to avoid under or overestimation of nutrient use.

Specific survey questions should address the following:

\begin{itemize}
\tightlist
\item
  What is the amount of nutrient input, such as nitrogen (N), phosphorus
  (P), or potassium (K), added to the farm system through fertilization
  or other means per year?
\item
  What proportion of these nutrients come from each source (chemical,
  organic plant, manure)?
\end{itemize}

\section{Guidance on data entry and
reporting}\label{guidance-on-data-entry-and-reporting-3}

Uniform data entry and conversion of inputs to kg ha\textsuperscript{-1}
yr\textsuperscript{-1}.

\section{Indicator interpretation and threshold
setting}\label{indicator-interpretation-and-threshold-setting-3}

\section{Limitations}\label{limitations-3}

\section{References}\label{references-3}

\begin{itemize}
\tightlist
\item
  Dobermann, (2007)\textsuperscript{21}
\item
  European Union Nitrogen Use Efficiency (2016)\textsuperscript{22}
\end{itemize}

\part{3. Social KPIs}

\chapter{3.1 Diet quality}\label{diet-quality}

\begin{longtable}[]{@{}
  >{\raggedright\arraybackslash}p{(\linewidth - 2\tabcolsep) * \real{0.1806}}
  >{\raggedright\arraybackslash}p{(\linewidth - 2\tabcolsep) * \real{0.8056}}@{}}
\toprule\noalign{}
\begin{minipage}[b]{\linewidth}\raggedright
\textbf{I ndicator}
\end{minipage} & \begin{minipage}[b]{\linewidth}\raggedright
\textbf{Diet quality}
\end{minipage} \\
\midrule\noalign{}
\endhead
\bottomrule\noalign{}
\endlastfoot
Theme & Nutrition \\
Purpose of indicator & Eradicating food insecurity and malnutrition is a
sustainable development goal (SDG2). Nutritious diets are a prerequisite
to achieving this goal. \\
Data source & Household and individual survey using the \textbf{Diet
Quality Questionnaire (DQQ)}, developed by the Global Diet Quality
Project \\
M easurement & Diet quality from 24-hour food consumption recalls,
calculated from the number of food groups consumed per individual and
household, from 29 food groups. \\
M easurement units & \begin{minipage}[t]{\linewidth}\raggedright
Raw measurement: Y/N for consumption of each food group together with
(optional) food group source (own production, purchased, borrowed, food
aid, other).

Final measurement:

\begin{itemize}
\item
  Household Food Group Diversity Score (FGDS) (0-10) (\textbf{this
  measurement is the default KPI})
\item
  Minimum Dietary Diversity for Women of Reproductive Age (MDD-W) (0/1)
\item
  All-5 recommended food groups consumed (0/1)
\item
  Non-Communicable Disease (NDC) Protect (0-9), which is an indicator of
  dietary factors that protect against NDCs according to the World
  Health Organisation (WHO)
\item
  NDC-Risk (0-9), which is an indicator of dietary factors increasing
  the risk of NCDs according to WHO
\item
  Global Dietary Recommendations (GDR) Score (0 to 18), which indicates
  level of adherence to global dietary recommendations for a healthy
  diet according to WHO
\item
  (Optional) Proportion of food that is produced on-farm
\end{itemize}
\end{minipage} \\
Guidance on m easurement & \begin{minipage}[t]{\linewidth}\raggedright
The \textbf{Diet Quality Questionnaire (DQQ)} is a standardized tool to
collect indicators of \textbf{\emph{dietary adequacy}}, including the
minimum dietary diversity for women (MDD-W) indicator, and All-5, as
well as indicators of \textbf{\emph{protection of health against
noncommunicable diseases (NCDs)}}, including NCD-Protect, NCD-Risk, and
the global dietary recommendations score (GDR). It is a low-burden tool
for collecting valid, comparable food group consumption data in
populations.

The DQQ is two-pages long and includes 29 multiple choice questions
corresponding to the 29 DQQ food groups. Each question asks the
respondent if they consumed a specific DQQ food groups with yes/no
responses. The questionnaire takes approximately 5 minutes to complete
per individual.

The survey should be completed for the adult household members,
including men and women. Researchers can optionally add one column
asking, for each food group, the source: own production, purchased,
borrowed, food aid, other. This will allow calculation of the proportion
of consumed food that is from each source.

Note that the timing of survey affects results. Try to collect data from
all farms at a period when food supplies are adequate, rather than when
some farms are more likely to suffer shortages (e.g.~end of dry season,
or after a drought season).

\textbf{DQQ Food Groups:}

1. Foods made from grains

2. Whole grains

3. White roots, tubers, and plantains

4. Pulses

5. Vitamin A-rich orange vegetables

6. Dark green leafy vegetables

7. Other vegetables

8. Vitamin A-rich fruits

9. Citrus

10. Other fruits

11. Baked / grain-based sweets

12. Other sweets

13. Eggs

14. Cheese

15. Yogurt

16. Processed meats

17. Unprocessed red meat (ruminant)

18. Unprocessed red meat (non-ruminant)

19. Poultry

20. Fish and seafood

21. Nuts and seeds

22. Packaged ultra-processed salty snacks

23. Instant noodles

24. Deep fried foods

25. Fluid milk

26. Sweet tea / coffee / cocoa

27. Fruit juice and fruit-flavored drinks

28. Sugar-sweetened beverages (soft drinks, energy drinks, sports
drinks)

29. Fast food

\textbf{Examples of how the final scores are calculated}:

\begin{figure}[H]

{\centering \includegraphics[width=4.7607in,height=2.47002in]{kpis/social/media/image4.png}

}

\caption{A table with numbers and text Description automatically g
enerated}

\end{figure}%

\begin{figure}[H]

{\centering \includegraphics[width=5.16225in,height=0.99242in]{kpis/social/media/image5.png}

}

\caption{A white sheet with black text Description automatically g
enerated}

\end{figure}%
\end{minipage} \\
Guidance on data entry and reporting & \\
Indicator inte rpretation and threshold setting & \\
L imitations & \\
References & \begin{minipage}[t]{\linewidth}\raggedright
Global Diet Quality Project resources\textsuperscript{23}:

\begin{itemize}
\item
  \url{https://www.globaldietquality.org/dqq}
\item
  \href{https://drive.google.com/fil\%20e/d/1glD4SAdYKKCfTE2fWfr0VsZwpYC3L-3G/view?usp=sharing}{DQQ
  Food Group Descriptions}
\item
  \href{https://docs.google.com/document/d/1RGwNit\%20itbYqSrGpUXhQs7mHWgLIgOxAlN5N0AOZN3qc/edit?usp=sharing}{Diet
  Quality Questionnaire (DQQ) Indicator Guide}
\end{itemize}

Diet Quality Questionnaires are adapted to each country, downloadable
from the Global Diet Quality Project website and translated into many
languages.
\end{minipage} \\
\end{longtable}

\chapter{3.2 Land tenure security}\label{land-tenure-security}

\begin{longtable}[]{@{}
  >{\raggedright\arraybackslash}p{(\linewidth - 2\tabcolsep) * \real{0.2222}}
  >{\raggedright\arraybackslash}p{(\linewidth - 2\tabcolsep) * \real{0.7639}}@{}}
\toprule\noalign{}
\begin{minipage}[b]{\linewidth}\raggedright
\textbf{Indicator}
\end{minipage} & \begin{minipage}[b]{\linewidth}\raggedright
\textbf{Land tenure security}
\end{minipage} \\
\midrule\noalign{}
\endhead
\bottomrule\noalign{}
\endlastfoot
Theme & Land tenure \\
Purpose of indicator & Land access is closely linked to use and control
over productive resources. Secure land tenure is therefore a key
determinant of individual agency, income and wellbeing. This indicator
is aligned with SDG 1.4.2 and 5.a.1. \\
Data source & Parcels (most preferable) or household/farm (most
possible) survey using a simplified version of the harmonized
questionnaire proposed by FAO et al. (2019)\textsuperscript{24}.
Additionally, GPS delineation of parcels is strongly advised wherever
feasible, as evidence points to systematic bias in farmer estimates of
land area\textsuperscript{25}. \\
Measurement & ((Area of agricultural land used by the household for
which at least one member of the household has their name on a legally
recognized document OR has the right to sell OR has the right to
bequeath)/(Total land area used by the household))*100 \\
Measurement units & \% \\
Guidance on measurement & \begin{minipage}[t]{\linewidth}\raggedright
The indicator can be calculated with data collected using FAO's
harmonized questionnaire, which is a standardized and succinct survey
instrument designed to collect the essential data for computation of
both 1.4.2\textsuperscript{a} and 5.A.1\textsuperscript{b} Sustainable
Development Goals (SDGs) indicators simultaneously.

The methodology proposed 5 versions of the questionnaire (see next
figure) depending on: (i) the respondent selection approach, (ii)
whether or not the survey collects data at the parcel-level and (iii)
whether or not the existing survey includes a parcel roster. Each
version of the questionnaire has approximated 10 questions and assumes
that a separate household member roster exists and that the sex of each
household member is captured in that roster.

\begin{figure}[H]

{\centering \includegraphics[width=3.89583in,height=2.84631in]{kpis/social/media/image6.png}

}

\caption{A diagram of a company Description automatically generated with
medium confid ence}

\end{figure}%

Where possible, it is advised to have parcel level data and implement a
self-respondent approach (i.e., where individuals respond for themselves
only, not on behalf of others) (version 2 with 7 questions). Where it is
not, land tenure data should be collected at the household/farm level,
using the version 3 (13 questions) of the questionnaire.

A sub-indicator can be calculated as: ((Number of adult women in
agricultural households whose name is on a legally recognized document
OR who have the right to sell OR bequeath agricultural land)/(Total
adult individuals in agricultural households whose name is on a legally
recognized document OR who have the right to sell OR bequeath
agricultural land))*100.

Note, for pastoralists, land tenure security may be better measured in
terms of mobility, i.e. proportion of grazing lands over which the
farmer has unrestricted rights to graze livestock.

\ul{Guidance of how to apply the questionnaire:}

- The survey should be completed for the adult household member (18
years old or older).

- Because a self-respondent approach is going to be used and only some
(or one) randomly selected household members are interviewed about their

tenure rights, only the selected members are part of the reference
population (i.e., denominator in the indicator and sub-indicator
formulas).

- The country-specific questionnaire should include: i) a comprehensive
list of all tenure types applicable to the country; ii) a comprehensive
list of land tenure-related documents, specifying which ones the
government considers as legally recognized; iii) \textbf{i}mages of the
documents considered legally recognized; iv) a context-specific
definition of alienation rights.

\begin{figure}[H]

{\centering \includegraphics[width=4.08333in,height=4.83333in]{kpis/social/media/image7.png}

}

\caption{A screenshot of a document Description automatically gene
rated}

\end{figure}%

\begin{figure}[H]

{\centering \includegraphics[width=4.83333in,height=4.19792in]{kpis/social/media/image8.png}

}

\caption{A close up of a document Description automatically gene rated}

\end{figure}%

\begin{figure}[H]

{\centering \includegraphics[width=4.83333in,height=2.78125in]{kpis/social/media/image9.png}

}

\caption{A screenshot of a computer Description automatically gene
rated}

\end{figure}%

\textsuperscript{a} Proportion of total adult population with secure
tenure rights to land, with (a) legally recognized documentation; and
(b) who perceive their rights to land as secure, by sex and by type of
tenure.

\textsuperscript{b} (a) proportion of total agricultural population with
ownership or secure rights over agricultural land by sex; and (b) share
of women among owners or rights bearers of agricultural land, by type of
tenure.
\end{minipage} \\
Guidance on data entry and reporting & Check that the total land area
accessed by the household matches the sum of the land areas with
different access types. \\
Indicator i nterpretation and threshold setting & \\
Limitations & \\
References & FAO; The World Bank; UN-Habitat. 2019. Measuring
Individuals' Rights to Land: An Integrated Approach to Data Collection
for SDG Indicators 1.4.2 and 5.a.1. Washington, DC.
72p.\textsuperscript{24} \\
\end{longtable}

\chapter{3.3 Farmer agency}\label{farmer-agency}

\begin{longtable}[]{@{}
  >{\raggedright\arraybackslash}p{(\linewidth - 2\tabcolsep) * \real{0.1667}}
  >{\raggedright\arraybackslash}p{(\linewidth - 2\tabcolsep) * \real{0.8194}}@{}}
\toprule\noalign{}
\begin{minipage}[b]{\linewidth}\raggedright
\textbf{In dicator}
\end{minipage} & \begin{minipage}[b]{\linewidth}\raggedright
\textbf{Farmer agency}
\end{minipage} \\
\midrule\noalign{}
\endhead
\bottomrule\noalign{}
\endlastfoot
Theme & Social justice \\
Purpose of indicator & Individual agency to decide on what to consume,
produce, and trade, and how to produce it, is a basic human right for
farmers. Monitoring farmer agency provides an indication of the power
dynamics influencing major decisions about the farm and food system, and
level of social equality (SDG 10, 16). \\
Data source & Household surveys \\
Me asurement & Ladder of Power and Freedom, from GENNOVATE \\
Me asurement units & Score of 1 to 5, and quotes and insights from
analysing qualitative data, for the household and each individual
surveyed \\
Guidance on me asurement & \begin{minipage}[t]{\linewidth}\raggedright
The GENNOVATE Ladder of Power and Freedom is a qualitative data
collection tool aimed at understanding local men and women farmers' own
perceptions of agency and key factors determining levels of agency.
Here, agency refers to `the capacity to make important decisions in
one's life and act upon them' \textsuperscript{26}.

The Ladder provides comparative qualitative measures of agency that
remain contextually grounded, and can be converted to a single
quantitative agency score (but see notes on effectively applying the
ladder approach). Data are collected through an individual interview
using the ladder visual (see inset), which is a scoreboard showing
different levels of agency. Data can also be collected at the community
level through a focus group.

\begin{figure}[H]

{\centering \includegraphics[width=5.21616in,height=4.23256in]{kpis/social/media/image10.png}

}

\caption{A close-up of a questionnaire Description automatically g
enerated}

\end{figure}%

\emph{Source:} \textsuperscript{26}

The ladder is applied as follows. The researcher identifies a meaningful
historical benchmark, usually 10 years in the past, and provides (or
creates) a visual of the five-step ladder on a sheet of paper. The
researcher introduces the purpose of the exercise. They will need to
explain the five-step ladder, and that step 5 indicates great power and
freedom while step 1 indicates very little power and freedom to make
important decisions.

For example: \emph{Today we want to talk to you about how much power and
freedom {[}wo{]}men in this household have to control the household food
system -- that is, how and what kind of food is produced, where it goes,
and how and what kind of food people consume. Please imagine a five-step
ladder {[}show figure{]}, where at the bottom, on the first step, stand
the individual {[}wo{]}men of this community with little capacity to
make their own decisions about food production and consumption. These
{[}wo{]}men have little to say about the types and varieties of crops,
animals, trees, and animals they cultivate, the way they manage their
water, land, soil, animals, and wildlife, the markets available to them,
and the food source and food choices available to them. On the highest
step, the fifth, stand those who have great capacity to make important
decisions for themselves, including what they grow, how they manage
their land, soil, water, animals, and trees, and what food they buy.}

Request the individual to mark on the ladder the step they think
household members of their gender occupy.

Facilitate the discussion on reasons for the step identified. For
example:

\emph{Q. Why would you place {[}wo{]}men on this step? Would you like to
volunteer the reasons for your rating?}

\emph{Q. Would you position different types of {[}wo{]}men in different
places on this ladder? Older people, younger people, poorer people,
single people, etc.? Do women and men differ?}

Repeat the rankings and discussion to capture perceptions of power and
freedom a decade ago (or at the pre-defined historical benchmark). These
sticky notes / cards can fill in the other side of the ladder. These
discussions should also be probed deeply to elicit detailed explanations
for the trends in agency identified. The goal of the exercise is to
understand the factors that shape women's or men's conceptions of power
and freedom in their lives, and reasons for changes in these conceptions
over time, rather than obtaining exact measurements or absolute values
for agency.

Time permitting, facilitate a discussion about the changes the
participants would like to see and what would be needed to achieve their
vision. For example:

\emph{Q. Having talked through this together, are there changes you
would like to see to give {[}wo{]}men more power and freedom in the
future? What would be needed to achieve this? What are the major
obstacles?}

Record the whole interview and take good notes. During the discussion,
listen closely for the reasoning behind votes, root causes of high or
low agency, factors that contributed to change, and any differences
reported for different groups in the community. Make note of powerful
quotes that could be extracted later. Record how many votes are cast for
each ladder step for both past and present.

\textbf{Example responses for data collected at the community level:}

\begin{figure}[H]

{\centering \includegraphics[width=5.16279in,height=2.66965in]{kpis/social/media/image11.png}

}

\caption{A table with numbers and arrows Description automatically g
enerated}

\end{figure}%

\emph{Source:} \textsuperscript{26}

\textbf{Notes on effectively applying the ladder approach:}

\begin{itemize}
\item
  Researchers should ideally be the same sex as participants.
\item
  Individual surveys should be conducted in a private room without the
  presence of any members of the household or anyone of the opposite
  sex.
\item
  Since the tool requires respondents to recall circumstances 10 years
  ago, the minimum recommended respondent age is 25 years old. However
  the 10-year timeframe could be reduced if younger respondents are
  present.
\item
  To assist with recall, it is helpful to identify the year
  corresponding to 10 years ago (e.g.~2014) and an important event at
  that time.
\item
  Researchers need to have a good level of skill in qualitative data
  collection to collect meaningful data on agency using this (or any)
  method. For example, managing group dynamics in household (or focus
  group) surveys is important to ensure the perspectives of individuals
  with lower levels of agency are captured.
\item
  The researcher should probe deeply into individual explanations of
  ladder rankings to understand the reasons for these rankings.
\item
  Reducing agency to a single score (i.e., the mean step change at
  household level) risks undermining the relevance and quality of the
  data collected. Focusing on qualitative data collection and analysis
  is strongly recommended for a more meaningful and accurate
  understanding of individual and household agency.
\end{itemize}
\end{minipage} \\
Guidance on data entry and reporting & This indicator collects sensitive
data on human subjects. Researchers should obtain ethical approval from
a verified review board prior to data collection. Study participants
should receive assurances that their testimonies will be provided
anonymously and no identifying information shared in any future
publications. \\
Indicator inter pretation and threshold setting &
\begin{minipage}[t]{\linewidth}\raggedright
\begin{itemize}
\tightlist
\item
\end{itemize}
\end{minipage} \\
Li mitations & \begin{minipage}[t]{\linewidth}\raggedright
\begin{itemize}
\tightlist
\item
\end{itemize}
\end{minipage} \\
R eferences & \begin{minipage}[t]{\linewidth}\raggedright
\begin{itemize}
\item
  Petesch and Bullock (2018) \textsuperscript{26}
\item
  Tavenner and Crane (2022) \textsuperscript{27}
\end{itemize}
\end{minipage} \\
\end{longtable}

\chapter{3.4 Human well-being}\label{human-well-being}

\begin{longtable}[]{@{}
  >{\raggedright\arraybackslash}p{(\linewidth - 2\tabcolsep) * \real{0.2222}}
  >{\raggedright\arraybackslash}p{(\linewidth - 2\tabcolsep) * \real{0.7639}}@{}}
\toprule\noalign{}
\begin{minipage}[b]{\linewidth}\raggedright
Indicator
\end{minipage} & \begin{minipage}[b]{\linewidth}\raggedright
Personal Wellbeing Index
\end{minipage} \\
\midrule\noalign{}
\endhead
\bottomrule\noalign{}
\endlastfoot
Theme & Human well-being \\
Purpose of indicator & This indicator aligns with SDG 3 ``Ensure healthy
lives and promote well-being for all at all ages''. \\
Data source & Household survey \\
Measurement & OECD Personal Wellbeing Index evaluation questions \\
Measurement units & \begin{minipage}[t]{\linewidth}\raggedright
\phantomsection\label{63}{I would drop this indicator, which is tailored
to individuals in developed countries. Subjective wellbeing indicators
need to be anchored to some frame-of-references, which are other
individuals in the community to which the individual situation should be
compared.}\phantomsection\label{64}{I agree-\/-these questions seem to
be from the OECD tool-\/-which I assume are targeted towards wealthy
OECD countries with predominantly urban
population}\phantomsection\label{65}{instead consider Generalized
self-efficacy score used in Pro-WEAI. I cope text below from
\href{../customXml/item1.xml}{https://weai.ifpri.info/files/
2023/07/Pro-WEAI-Guide.pdf}\\
The NGSE is a cross-culturally validated scale to measure self-efficacy,
or a person's perception of their capabilities and ability to reach
their goals (Chen et al., 2001). We ask respondents to think about how
each statement relates to their life and indicate if they strongly
disagree, disagree, neither agree nor disagree, agree or strongly agree
with the statement. The four statements are: 2) When facing difficult
tasks, I am certain that I will accomplish them, 3) In general, I think
that I can obtain outcomes that are important to me, 6) I am 9 confident
that I can perform effectively on many different tasks, and 8) Even when
things are tough, I can perform quite well.5 For each statement that an
individual strongly disagrees with they are given a score of 1. If they
disagree, they are given a score of 2, neither agree nor disagree
receives a score of 3, agree receives a score of 4 and strongly agree
receives a score of 5. The scores for each of the 4 statements are then
summed to get the final score. For example, if an individual strongly
agrees with all 4 statements they will have a score of 20 (4x5). In
order to be adequate in self efficacy a respondent must have a New
General Self-Efficacy Scale score of at least 16. This is equivalent to
an average response across the statements of ``agree."}Likert scale
(0-10)\phantomsection\label{63}{\phantomsection\label{64}{\phantomsection\label{65}{}}}\strut
\end{minipage} \\
Guidance on measurement & \begin{minipage}[t]{\linewidth}\raggedright
Human well-being encompasses quality of life and the ability of people
and societies to contribute to the world with a sense of meaning and
purpose. It encompasses a broad range of factors, including material
well-being and emotional, physical, and social
health\textsuperscript{28}. While people's living conditions (e.g.,
housing, income, employment) are important contributors to well-being,
measuring living conditions (i.e.~material well-being) alone fails to
capture what people think and feel about their lives -- i.e., their
subjective well-being.

OECD's Personal Wellbeing Index is computed based on responses to the
following questions\textsuperscript{29}:

\begin{enumerate}
\def\labelenumi{\arabic{enumi}.}
\item
  Thinking about your own life and personal circumstances, how satisfied
  are you with your life as a whole? {[}0-10{]}
\item
  How satisfied are you with your standard of living? {[}Standard of
  Living{]} {[}0-10{]}
\item
  How satisfied are you with your health? {[}Personal Health{]}
  {[}0-10{]}
\item
  How satisfied are you with what you are achieving in life?
  {[}Achieving in Life{]} {[}0-10{]}
\item
  How satisfied are you with your personal relationships? {[}Personal
  Relationships{]} {[}0-10{]}
\item
  How satisfied are you with how safe you feel? {[}Personal Safety{]}
  {[}0-10{]}
\item
  How satisfied are you with feeling part of your community?
  {[}Community-Connectedness{]} {[}0-10{]}
\item
  How satisfied are you with your future security? {[}Future Security{]}
  {[}0-10{]}
\item
  How satisfied are you with the amount of time you have to do the
  things that you like doing? {[}0-10{]}
\item
  How satisfied are you with the quality of your local environment?
  {[}0-10{]}
\item
  For respondents who are employed only: How satisfied are you with your
  job? {[}0-10{]}
\end{enumerate}
\end{minipage} \\
Guidance on data entry and reporting & Key issues include the need to be
careful with question wording, use consistent scales, and appropriate
question ordering.

1. Question wording

``Because question wording matters to how people respond to questions on
subjective well-being, the translation of questions is of high
importance. Potential issues arising from translation cannot be
eliminated, but they can be managed through an effective translation
process. A robust translation process, including back translation, is
therefore essential'' \textsuperscript{29}.

2. Scale design

``Variation in response formats can affect data quality and
comparability. In the case of evaluative measures, there is empirical
support for using a 0-10 point numerical scale, anchored by verbal
labels which represent conceptual absolutes (such as completely
satisfied/completely dissatisfied). On balance, it is preferable to
label scale interval points (between the anchors) with numerical, rather
than with verbal, labels'' \textsuperscript{29}.

3. Question order

\phantomsection\label{66}{Replies are strongly affected by individual
perceptions about own status compared to neighbors\textquotesingle{} or
friends\textquotesingle, so they are expression on a relative -not
absolute- condition.}Subjective well-being questions are influenced by
context \phantomsection\label{66}{}and thus need to be asked in the
right place so as not to be influenced by previous questions. {Also
cons}{ider randomizing the order in which the well-being questions are
asked to reduce ordering-bias.} \\
Indicator i nterpretation and threshold setting & \\
Limitations & \\
References & WHO (2024), Promoting well-being, \textless h
ttps://www.who.int/activities/promoting-well-being\textgreater, accessed
1 July 2024\textsuperscript{28}

OECD (2013), OECD Guidelines on Measuring Subjective Well-being, OECD
Publishing.
\url{http://dx.doi.org/10.1787/9789264191655-en}\textsuperscript{29}

International Wellbeing Group (2006), Personal Wellbeing Index, 4th
edition, Melbourne, Australian Centre on Quality of Life, Deakin
University.\textsuperscript{30} \\
\end{longtable}

\part{4. Economic KPIs}

\chapter{4.1 Agricultural productivity}\label{agricultural-productivity}

\begin{longtable}[]{@{}
  >{\raggedright\arraybackslash}p{(\linewidth - 2\tabcolsep) * \real{0.2222}}
  >{\raggedright\arraybackslash}p{(\linewidth - 2\tabcolsep) * \real{0.7639}}@{}}
\toprule\noalign{}
\begin{minipage}[b]{\linewidth}\raggedright
\textbf{Indicator}
\end{minipage} & \begin{minipage}[b]{\linewidth}\raggedright
\textbf{Crop and livestock productivity}
\end{minipage} \\
\midrule\noalign{}
\endhead
\bottomrule\noalign{}
\endlastfoot
Theme & Agricultural productivity \\
Purpose of indicator & Ensuring productivity and closing yield gaps is a
primary goal in agricultural systems and vital for achieving global food
security (SDG2), and securing livelihoods (SDG1). \\
Data source & Household survey (farmer recall) \\
Measurement & For crops: kg/ha (dry weight) of
\phantomsection\label{75}{Why "harvestable" and not
"harvested"?}\phantomsection\label{76}{totally agreed - we have to be
specific}\phantomsection\label{77}{Also is "dry-weight" needed? For
example: maize etc consumed as corn on the cob, dry weight is not
relevant, but maize used as feed might use
dry-weight?}\phantomsection\label{78}{I think harvestable is to qualify
that this is non-spoilt/infected crop. Perhaps a better qualifier is
consumable/usable/sellable crop?}harvestable
\phantomsection\label{75}{\phantomsection\label{76}{\phantomsection\label{77}{\phantomsection\label{78}{}}}}produce;
land equivalent ratio for intercropped and agroforestry systems (0 to
infinity)

For livestock: \#\# animals slaughtered for consumption/sale/gifts, kg
meat/animal, litres (milk)/animal {eggs/chicken, manure/animal}{,
wool/animal,} {}{skin/animal} \\
Measurement units & Kg/ha; kg/animal; l/animal; land equivalent area \\
Guidance on measurement & \begin{minipage}[t]{\linewidth}\raggedright
\textbf{\ul{Crop productivity through farmer recall}}

Harvestable production per unit land area (kg/ha) for at least
\phantomsection\label{83}{Why only three main crops? And main based on
land area, sales, quantity harvested,
or...?}\phantomsection\label{84}{agree - each crop grown by each
household has some importance - why then limit the farmer - if they can
recall info for all crops, we should allow
it.}\phantomsection\label{85}{It says minimum 3, not maximum. We have to
make tradeoffs based on length of survey and respondent fatigue.\\
Farmers usually know what their "main crops" are when you ask
them-\/-usually based on economic importance to the HHs, not necessarily
quantity or land (though most times it is correlated)}the three main
crops\phantomsection\label{83}{\phantomsection\label{84}{\phantomsection\label{85}{}}},
based on farmer recall for all produce harvested over the past 12
months. In the case of intercropped or agroforestry systems, collect
data for all crops present in the plot and compute the land equivalent
ratio (LER).

\textbf{\ul{Livestock productivity}}

Production of animal products per unit livestock, for at least the three
main livestock species (including fish), based on farmer recall of all
products consumed, sold or gifted over the last 12 months.

Desirable additional data:

\begin{itemize}
\tightlist
\item
  Population growth rate: ratio between the production in number of
  animals and the initial population size
\item
  Total annual livestock production = (stock variation +offtake-Intake)/
  Initial size
\item
  Production of milk : av. production per female (/day)* number of
  reproductive female* milking period
\item
  Skin \& hides (from offtake)
\item
  Wool : total production per year
\item
  Manure : total production (can be collected in declarative)
\item
  Productivity indicators: 1) Number of new sub-adults produced per
  reproductive female per year 2) Number of new adults produced per
  reproductive female and year
\end{itemize}

\textbf{\ul{Crop productivity - detailed measurements (beyond the scope
of HOLPA)}}

For crops see FAO (2017): https://www.fao.org/3/ca6514en/ca6514en.pdf.

For livestock see FAO (2018):
https://www.fao.org/3/ca6400en/ca6400en.pdf.

Summary for crop yield: Yields of all crops should be measured by
randomly demarcating at least three same-size plots in a field,
harvesting all produce, then drying the produce and weighing it to
determine its dry weight. If possible take into account harvest losses
(proportion of diseased, rotten or otherwise unusable produce).
Alternatively, yields can be measured based on farmer recall
particularly where farmers keep good records, but note that in many
cases this method is inaccurate. Alternatively, farmers can be asked to
keep crop diaries and record all harvested produce over the cropping
season, which produced reliable results for the LSMS but is unsuitable
for illiterate farmers. See guide for further details on calculating
yields in mixed cropping systems.

Crop and pasture areas should be estimated by walking along the
perimeter of the farmed areas with a GPS (preferred method) or using the
rope-and-compass method (see
https://www.fao.org/3/ca6514en/ca6514en.pdf). Areas of additional plots
that are far away or difficult to access can be estimated by farmers,
and verified on Google Earth imagery if possible. Asking farmers to
estimate the area of their fields can work well for small parcel sizes
in some contexts, but is often unreliable e.g.~when farmers are not
familiar with the units of measurement or if farmers have incentives to
misreport crop area (see
https://www.fao.org/3/ca6514en/ca6514en.pdf).\strut
\end{minipage} \\
Guidance on data entry and reporting & \\
Indicator i nterpretation and threshold setting & \\
Limitations & \\
References & Dutilly C., Alary V., Bonnet P., Lesnoff M., Fandamu P., De
Haan Cees, 2020. Multi-scale assessment of the livestock sector for
policy design in Zambia. Journal of Policy Modeling 42 (2): 401-418.
\url{https://doi.org/10.1016/j.jpolmod.2019.07.004}\textsuperscript{31}

Alary, V., Lasseur, J., Frija, A., \& Gautier, D. (2022). Assessing the
sustainability of livestock socio-ecosystems in the drylands through a
set of indicators. Agricultural Systems, 198,
103389.\textsuperscript{32} \\
\end{longtable}

\chapter{4.2 Labour productivity}\label{labour-productivity}

\begin{longtable}[]{@{}
  >{\raggedright\arraybackslash}p{(\linewidth - 2\tabcolsep) * \real{0.2222}}
  >{\raggedright\arraybackslash}p{(\linewidth - 2\tabcolsep) * \real{0.7639}}@{}}
\toprule\noalign{}
\begin{minipage}[b]{\linewidth}\raggedright
Indicator
\end{minipage} & \begin{minipage}[b]{\linewidth}\raggedright
\phantomsection\label{88}{This is labor intensity, not
productivity.}Labour inputs per unit agricultural land
are\phantomsection\label{88}{}a
\end{minipage} \\
\midrule\noalign{}
\endhead
\bottomrule\noalign{}
\endlastfoot
Theme & Labour productivity \\
Purpose of indicator & Labour inputs required for agricultural
production is one of the primary factors determining which practices are
used, how profitable they are, and the social sustainability of the
production system. Distinguishing which productive systems promote
sustainable employment contributes to SDG 2 and SDG8. \\
Data source & Farmer recall during household survey \\
Measurement & Person hours per year, disaggregated by gender, season,
and productive activity \\
Measurement units & \phantomsection\label{89}{This is amount of labor,
it\textquotesingle s not intensity or productivity.}Person hours per
year {per} {hectare} {o}{f cultivated
land}\phantomsection\label{89}{} \\
Guidance on measurement & \begin{minipage}[t]{\linewidth}\raggedright
This indicator can be computed by collecting data on number of labour
hours per year per hectare of land actively used for agricultural
purposes. Labour hours should be disaggregated by:

\begin{itemize}
\item
  gender
\item
  age
\item
  type of worker: hired or not hired, permanent or seasonal
\item
  type of activity: routine (e.g.~manure collection, milking) or
  seasonal work (e.g. seeding, harvesting, pruning).
\end{itemize}
\end{minipage} \\
Guidance on data entry and reporting & \\
Indicator i nterpretation and threshold setting & \\
Limitations & \\
References & FAO (2010) \textsuperscript{33} \\
\end{longtable}

\chapter{4.3 Total Income and Income
Stability}\label{total-income-and-income-stability}

Content for this KPI will be added here.

\chapter{4.4 Climate resilience}\label{climate-resilience}

\begin{longtable}[]{@{}
  >{\raggedright\arraybackslash}p{(\linewidth - 2\tabcolsep) * \real{0.2222}}
  >{\raggedright\arraybackslash}p{(\linewidth - 2\tabcolsep) * \real{0.7639}}@{}}
\toprule\noalign{}
\begin{minipage}[b]{\linewidth}\raggedright
Indicator
\end{minipage} & \begin{minipage}[b]{\linewidth}\raggedright
Resilience score
\end{minipage} \\
\midrule\noalign{}
\endhead
\bottomrule\noalign{}
\endlastfoot
Theme & \begin{minipage}[t]{\linewidth}\raggedright
\phantomsection\label{136}{World Bank: resilient futures project\\
Maybe have indicators that we can calculate using the data we
collected?\\
Carlo also has some ideas on simple outcome indicators we might be able
to use}\phantomsection\label{137}{Comparison paper:
\href{../customXml/item2.xml}{http
s://www.sciencedirect.com/science/article/pii/S03043
8782200044X}}Climate
resilience\phantomsection\label{136}{\phantomsection\label{137}{}}\strut
\end{minipage} \\
Purpose of indicator & Climate change is affecting food systems around
the world, with temperatures and precipitation anomalies and extreme
events leading to crop/livestock losses, altering species ranges,
creating agricultural water stress, altering working conditions and
triggering human migration. Building farm and food system resilience to
climate change and other major stressors (e.g.~pandemics, war) is
fundamental to prevent food and livelihood insecurity (SDG 1, 2), and
avoidable loss of life (SDG 3, 15). \\
Data source & Household survey \\
Measurement & Resilience score adapted from RIMA \\
Measurement units & \phantomsection\label{138}{The RIMA methodology has
been subject to severe criticism and scrutiny. The list of component
variables is completely arbitrary and the final RIMA value is not a
measure of realized resilience, but rather a measure of wellbeing that
might or might not be correlated to actual household resilience to
shocks, The four pillars represent a grouping of variables proxies of
welfare and the resulting indicator is again a welfare proxy, which has
little conceptual relationship with resilience. Of course, household
with higher values in the component variables could likely be associated
to higher resilience to weather shocks, but only through their higher
welfare. It could very well be called a proxy indicator of
multidimensional poverty.}Score \phantomsection\label{139}{this scale is
too wide}0-20\phantomsection\label{138}{\phantomsection\label{139}{}} \\
Guidance on measurement & \begin{minipage}[t]{\linewidth}\raggedright
Resilience is the capacity of socioeconomic systems (e.g., households)
to withstand shocks through adaptation and transformation.

FAO's Resilience Index Measurement and Analysis (RIMA) is one way to
measure resilience. FAO computes a Resilience Capacity Index (RCI) based
on data collected across four pillars:

\begin{enumerate}
\def\labelenumi{\arabic{enumi}.}
\item
  Access to basic services (ABS)
\item
  Assets (AST)
\item
  Social safety nets (SSN)
\item
  Adaptive capacity (AC)
\end{enumerate}

FAO calculates the RCI using structural equation modelling (SEM),
combining Confirmatory Factor Analysis and Path Analysis. There are two
steps: 1) estimate resilience of each pillar through Factor Analysis
(FA), 2) calculate a resilience index from the pillars, using Ordinary
Least Squares with the outcome variable being a food security
indicator(s) such as dietary diversity score. The benefit of this
approach is that is possible to quantify the contribution of each pillar
to the overall resilience. The limitation is that the RCI depends
heavily on which indicator is used as the outcome variable which does
not capture the variation across the four pillars.

For HOLPA, we adapt the FAO approach to directly combine the data from
the four pillars into a measure of resilience scaled from 0-20. This can
be done by assigning scores of 0-5 to all responses, then taking the
mean within each pillar, and summing scores across the four pillars
according to the following table:

\begin{longtable}[]{@{}
  >{\raggedright\arraybackslash}p{(\linewidth - 4\tabcolsep) * \real{0.0694}}
  >{\raggedright\arraybackslash}p{(\linewidth - 4\tabcolsep) * \real{0.3750}}
  >{\raggedright\arraybackslash}p{(\linewidth - 4\tabcolsep) * \real{0.2222}}@{}}
\toprule\noalign{}
\begin{minipage}[b]{\linewidth}\raggedright
\textbf{ Pi ll ar }
\end{minipage} & \begin{minipage}[b]{\linewidth}\raggedright
\textbf{Question}
\end{minipage} & \begin{minipage}[b]{\linewidth}\raggedright
\textbf{Score assigned to response}
\end{minipage} \\
\midrule\noalign{}
\endhead
\bottomrule\noalign{}
\endlastfoot
A BS & Does your farm/household have the following services: 1) Drinking
water piped to the household, 2) Toilet with piped sewer or septic tank,
3) Electrical energy, 4) Waste (rubbish) collection service, 5) Mobile
phone reception, 6) Internet. {[}Count services{]} & 5: 3 or more

2.5: 1-2

0: None \\
A BS & How far (one way) is the household from the closest accessible
and functioning service in minutes (walking): 1) Freshwater source, 2)
primary school, 3) public hospital or health facility, 4) livestock
market 5) agricultural/crops market, 6) public means of transport, 7)
well-maintained road suitable for cars. & For each service:

5: \textless5 mins

2.5: 5-20 mins

0: \textgreater{} 20 mins \\
A ST & How many of the following assets do household members own: 1)
Car, 2) Motorbike, 3) Bicycle, 4) Gas or electric cooker, 5) Mobile
phone, 6) Ox-plough, 7) Tractor, 8) Machete. {[}Count assets{]} & 5: 5
or more distinct asset types

2.5: 2-4

0: 0-1 \\
S SN & On average, how many school meals have the children living in the
household received per week, during the last month of attending school?
& 5: 5 meal

4: 4 meal

3: 3 meal

2: 2 meal

1: 1 meal

0: 0 meals \\
S SN & Would any of the following individuals or organisations provide
your household support in case of need: 1) local farmer associations or
cooperatives, 2) other local associations (e.g. women support groups,
youth groups), 3) individuals from my community, 4) individuals from a
different community, 5) community leaders, 6) NGOs 8) local government,
9) national government, 10) other (please specify). {[}Count sources{]}
& 5: 3 or more

2.5: 1-2

0: None \\
AC & Can the head of household read and write in any language? & 5: Can
read and write

2.5: Can read or write

0: Cannot read or write, or prefer not to say \\
AC & On average, how many years have the household members of working
age (\textgreater14 and \textless65 years old) attended formal school? &
5: ≥7 (Secondary or higher)

2.5:\textless7 (Primary only)

0:0 (None) \\
AC & Has any household member received training in the last 12 months? &
5: Y

0: N \\
AC & In the past 12 months, what were all the household's sources of
income: : 1) Agriculture, animal breeding, fishing, 2) Other family
business, 3) Government salary, 4) private sector salary, 5) Cash
transfers (e.g. transfers from absent family members, donations), 6)
Other (please specify) {[}count sources{]} & 5: 3 or more sources

2.5: 2 sources

0: 1 source \\
AC & In the past 7 days, what percentage of the
household\textquotesingle s income was used for buying food & 5: 0-25\%

3.33: 26-50\%

1.66: 51-75\%

0: 76-100\% \\
AC & In the past 5 years, has your farm/household had access to credit
to support investment in your farming business. & 5:Y

0:N, or I don't know \\
AC & How do you feel about your farm/household\textquotesingle s ability
to meet credit repayments now and until the loans are fully repaid? & 5:
Full repayments

3.33: Uncertain future repayments

1.66: Uncertain current repayments

0: Uncertain current and future repayments \\
AC & Do you have access to crop insurance? If yes, what is covered, e.g.
losses from weather events, pest outbreaks market shocks? & 5:Y

0:N \\
AC & Do you have access to subsidies for agricultural inputs? If yes,
which inputs & 5:Y

0:N \\
\end{longtable}

Additionally, HOLPA collects data on household exposure to shocks over
the last 12 months and any coping mechanisms employed.
\end{minipage} \\
Guidance on data entry and reporting & \\
Indicator i nterpretation and threshold setting & \\
Limitations & \\
References & FAO (2016) RIMA\textsuperscript{34} \\
\end{longtable}

\part{5. Complementary Indicators}

\chapter{5.1 Bat, Bird, and Insect
Diversity}\label{bat-bird-and-insect-diversity}

Content for this complementary indicator will be added here.

\chapter{5.2 Pollinator Diversity}\label{pollinator-diversity}

Content for this complementary indicator will be added here.

\chapter{5.3 Soil Health (SOCLA)}\label{soil-health-socla}

Content for this complementary indicator will be added here.

\chapter{5.4 Nutrient Use Efficiency}\label{nutrient-use-efficiency}

Content for this complementary indicator will be added here.

\chapter{5.5 Nutrient Balance}\label{nutrient-balance}

Content for this complementary indicator will be added here.

\chapter{5.6 Greenhouse Gas Emissions}\label{greenhouse-gas-emissions}

Content for this complementary indicator will be added here.

\chapter{5.7 Enteric Methane Emission
Intensity}\label{enteric-methane-emission-intensity}

Content for this complementary indicator will be added here.

\bookmarksetup{startatroot}

\chapter{Agroecological Indicators}\label{agroecological-indicators}

Content for the agroecological indicators chapter will be added here.

\bookmarksetup{startatroot}

\chapter*{References}\label{references-4}
\addcontentsline{toc}{chapter}{References}

\markboth{References}{References}

\phantomsection\label{refs}




\end{document}
